% ===========================================================================
%  comune/colori.tex — la tavolozza
% ===========================================================================
%
%  Sono i colori del PRODOTTO, non colori scelti per la documentazione: gli
%  stessi valori che stanno in dashboard/static/css/ui-v2.css e nella landing.
%  Documentare un prodotto con i colori di qualcun altro costringe chi legge a
%  tradurre due volte — una figura verde acqua qui e un pulsante verde acqua
%  nell'app sono la stessa cosa, e si deve vedere.
%
%  Chi cambia la tavolozza del prodotto cambia QUESTO file e nient'altro:
%  nessun altro file della documentazione nomina un colore per esteso.
% ===========================================================================
\RequirePackage{xcolor}

% --- i colori di base della landing -----------------------------------------
\definecolor{lyraInchiostro}{HTML}{111318}  % testo, titoli
\definecolor{lyraPannello}{HTML}{F5F5F2}    % fondo dei riquadri
\definecolor{lyraBordo}{HTML}{ECECEA}       % filetti e bordi
\definecolor{lyraAccento}{HTML}{2BB59A}     % verde acqua: l'unico accento
\definecolor{lyraPositivo}{HTML}{0F7A62}    % numeri in utile
\definecolor{lyraNegativo}{HTML}{9A3412}    % numeri in perdita

% --- i pastelli, usati per distinguere gli strati nelle figure --------------
\definecolor{lyraCielo}{HTML}{D9F1F7}
\definecolor{lyraMenta}{HTML}{DDF3E4}
\definecolor{lyraPesca}{HTML}{FFE4D4}
\definecolor{lyraLavanda}{HTML}{E6DEFF}
\definecolor{lyraCrema}{HTML}{FFF1C9}

% --- derivati, per bordi e testi sopra i pastelli ---------------------------
%  Si ricavano dai precedenti invece di essere scritti a mano: cambiando un
%  pastello, il suo bordo lo segue.
\colorlet{lyraCieloBordo}{lyraCielo!55!lyraInchiostro}
\colorlet{lyraMentaBordo}{lyraMenta!55!lyraInchiostro}
\colorlet{lyraPescaBordo}{lyraPesca!55!lyraInchiostro}
\colorlet{lyraLavandaBordo}{lyraLavanda!55!lyraInchiostro}
\colorlet{lyraCremaBordo}{lyraCrema!55!lyraInchiostro}

\colorlet{lyraGrigio}{lyraInchiostro!55}
\colorlet{lyraGrigioChiaro}{lyraInchiostro!28}
\colorlet{lyraAccentoScuro}{lyraAccento!70!lyraInchiostro}
\colorlet{lyraAccentoTenue}{lyraAccento!14!white}

% --- i colori del codice ----------------------------------------------------
\colorlet{codFondo}{lyraPannello}
\colorlet{codBordo}{lyraBordo}
\colorlet{codCommento}{lyraInchiostro!45}
\colorlet{codParola}{lyraAccentoScuro}
\colorlet{codStringa}{lyraNegativo!85}
\colorlet{codNumero}{lyraInchiostro!70}
\colorlet{codRiga}{lyraInchiostro!35}
